\subsection{A stochastic variance state visible in the curve}\label{sec:varianceexample}
The stochastic scale in Proposition~\ref{prop:scaledconstruction} is supplied by an additional diffusion. A different question is whether volatility can be determined by the curve's own coefficients while those coefficients stay in a fixed maturity family. The following construction answers this positively. It uses two Brownian factors and one scheduled maturity breakpoint, rather than the one-factor ordinal loading of \eqref{eq:ordinal}.

Fix a meeting date $\tau>0$ and a decay rate $\kappa>0$. On the usual augmentation of the natural filtration of independent Brownian motions $W^1,W^2$, start three coefficients $Z_1,Z_2,Z_4$ at zero and stop them at $\tau$. The subscripts identify their exponential decay multiples. Define
\begin{align}
v_t&=2+\tanh Z_2(t),&
Z_1(t)&=\frac{t\wedge\tau}{\kappa}+W^1_{t\wedge\tau},\nonumber\\
Z_2(t)&=\int_0^{t\wedge\tau}\frac{v_s-2}{2\kappa}\dd s
                  +\int_0^{t\wedge\tau}\sqrt{v_s}\dd W^2_s,&
Z_4(t)&=-\int_0^{t\wedge\tau}\frac{v_s}{2\kappa}\dd s.
\label{eq:varianceexample}
\end{align}
The variance $v$ lies strictly between one and three. For $0\le t\le T$, set
\begin{equation}\label{eq:threeexpcurve}
f(t,T)=\begin{cases}
0,&T<\tau,\\
Z_1(t)e^{-\kappa(T-\tau)}+Z_2(t)e^{-2\kappa(T-\tau)}
                   +Z_4(t)e^{-4\kappa(T-\tau)},&T\ge\tau.
\end{cases}
\end{equation}
This example is defined for all maturities; it may be restricted to any finite horizon. Each curve is bounded in maturity by $|Z_1(t)|+|Z_2(t)|+|Z_4(t)|$. This is not a uniform bound over states.

\begin{proposition}[Stochastic variance within a three-exponent family]\label{prop:varianceexample}
System~\eqref{eq:varianceexample} has a unique continuous strong solution. The curve~\eqref{eq:threeexpcurve} satisfies both consistency conditions, and $P(\cdot,T)/B$ is a positive true martingale for every finite maturity $T$.

For $t<\tau$, the instantaneous variance of $Z_2$ is $v_t$, and
\[
\dd[v]_t=(1-\tanh^2Z_2(t))^2v_t\dd t>0.
\]
For $t\le\tau$, the forward rates at $\tau$, $\tau+\log(2)/\kappa$, and $\tau+2\log(2)/\kappa$ determine all three coefficients and hence $v_t$. The short rate has the almost surely nonzero jump
\begin{equation}\label{eq:varianceexamplejump}
\Delta r_\tau=W^1_\tau+\int_0^\tau\sqrt{v_s}\dd W^2_s,
\end{equation}
although every fixed-maturity forward rate is continuous through $\tau$.
\end{proposition}

The mechanism is explicit. For $t<\tau\le T$ and $y=T-\tau$, the volatility rows are $e^{-\kappa y}$ and $\sqrt{v_t}e^{-2\kappa y}$. Their required drift is
\[
\frac{e^{-\kappa y}-e^{-2\kappa y}}{\kappa}
  +\frac{v_t}{2\kappa}(e^{-2\kappa y}-e^{-4\kappa y}),
\]
which is exactly the drift generated by \eqref{eq:varianceexample}. The second exponential is both part of the first factor's drift and a noisy curve coordinate; its own drift closes at the fourth exponential, whose coefficient has finite variation. Thus making a drift-generated coordinate stochastic need not force further and further maturity shapes. Appendix~\ref{app:varianceexample} proves existence, curve recovery, and the true-martingale assertion.

The fixed exponents and their doubling relations are compatible with the restrictions in \citet{filipovic2000exponential}; this construction does not contradict those results. Its purpose is to exhibit stochastic variance recoverable from the curve, rather than classify all possible variance laws. The variance is specified by $2+\tanh Z_2$, not by an independently imposed square-root SDE. There is no diffusion after $\tau$, and the short-rate jump comes from crossing the maturity breakpoint, not from new information arriving at the date.
